\guidepart{violet}{Work faster}{Work faster}
  {Fewer clicks, less typing, more getting on with it.}
  {\partitem{\iSplit}{Tabs and split panes}
   \partitem{\iKeyboard}{Terminal superpowers}
   \partitem{\iWand}{Never type it twice}
   \partitem{\iFolder}{Files: SFTP and drag-and-drop}
   \partitem{\iArrows}{Port forwards}
   \partitem{\iFilm}{Recordings}}

% =========================================================================================
\chapter{Tabs and split panes}\label{ch:tabs}
\chapmeta{4 minutes}{You'll juggle a dozen sessions without losing any of them}

\begin{tldr}
Right-click a tab for everything. \keys{\mod+Shift+D} splits right, \keys{\mod+Shift+E} splits down.
\keys{\mod+1} to \keys{\mod+9} jump straight to a tab.
\end{tldr}

\section{Tabs}

\begin{description}
  \item[Move] Drag tabs about, or right-click\then\ui{Move Left} / \ui{Move Right}.
  \item[Pin] Right-click\then\ui{Pin}. Pinned tabs stay at the front.
  \item[Rename] Right-click\then\ui{Rename}. ``prod DB, DON'T TOUCH'' is a perfectly valid tab name.
  \item[Duplicate] Opens another session to the same profile.
  \item[Reconnect] Does what it says.
  \item[Peek] Hover a tab for its name, host, state and when it connected.
  \item[Close] The \ui{×}, \keys{\mod+W}, or right-click\then\ui{Close}. Connected tabs ask first (you can turn that
        off in Settings).
  \item[Switch] \keys{\mod+Tab} and \keys{\mod+Shift+Tab}, or \keys{\mod+1} to \keys{\mod+9}.
\end{description}

Lots of tabs? Scroll the tab strip with your mouse wheel or trackpad. Quit with a dozen open and they're all back next
time: \ui{Restore tabs at launch} (on by default) reconnects saved profiles and reopens local shells, splits and all,
in the folder they were in.

\section{Split panes}

Want two sessions side by side? Split the focused pane right with \keys{\mod+Shift+D} or down with
\keys{\mod+Shift+E} (also in the \menu{Session} menu, the right-click menu and the palette). The new pane asks
what to open: \ui{Local Machine} or any of your terminal profiles.

\shot{split.png}{A fresh split. Pick what goes in the new pane.}

Drag a divider to resize. Click \ui{×} on a pane's header to close it and its neighbour takes the space. Closing
the tab closes every pane in it.

\begin{quickwin}
Split right, open \code{htop} on one side and tail a log on the other. Now you look like a hacker in a film.
\end{quickwin}

% =========================================================================================
\chapter{Terminal superpowers}\label{ch:terminal}
\chapmeta{4 minutes}{You'll copy, search and scroll like you've used it for years}

\begin{tldr}
Select to copy, \keys{\mod+F} to search, Ctrl/\ph{E1C4}-click links. Right-click for the rest.
\end{tldr}

\section{Keyboard and mouse}

\renewcommand\arraystretch{1.45}
\begin{tabularx}{\linewidth}{@{}>{\leavevmode\heavy\color{ink}}p{24mm}X@{}}
  \toprule
  Copy   & Select text, then \keys{\mod+C}. Nothing selected? \keys{Ctrl+C} interrupts the program as usual.\\
  Paste  & \keys{\mod+V}. Pastes are bracketed when the program asks for it, so editors won't auto-indent your paste
           into a staircase. Pasting several lines into a shell asks first (turn that off in Settings).\\
  Select & Drag for a range, double-click a word, triple-click a line, \kbd{Shift}-click to extend.\\
  Links  & Hold \kbd{Ctrl} (\kbd{\ph{E1C4}}) and click a URL to open it in your browser. Real hyperlinks that
           programs print (\code{ls --hyperlink}, compilers) get a dotted underline; hover to see where they really go.\\
  Scroll & Mouse wheel, or \keys{Shift+PageUp}, \keys{Shift+PageDown}, \keys{Shift+Home}, \keys{Shift+End}.\\
  \bottomrule
\end{tabularx}

\begin{sidequest}[Why does the scroll wheel move the cursor in less?]
Full-screen programs (\code{less}, \code{vim}, \code{tmux}, \code{byobu}) use the terminal's ``alternate screen'', where the
wheel sends arrow keys instead of scrolling. Programs with mouse support (\code{htop}, \code{vim} with
\code{mouse=a}, \code{tmux} mouse mode) get your clicks directly. That's normal terminal behaviour, not a bug.
\end{sidequest}

\section{The right-click menu}

\ui{Copy}, \ui{Paste}, \ui{Select All}, \ui{Find}, \ui{Previous Prompt} / \ui{Next Prompt}, \ui{Copy Last Command
Output}, \ui{Send Snippet}, \ui{Port Forwards} (SSH), \ui{Split Right}, \ui{Split Down}, \ui{Export Output} (saves the
whole scrollback as a text file), \ui{Clear} and \ui{Reset}.

\section{Shell integration}

If your shell marks its prompts (fish 4, Starship and Powerlevel10k do already), Nexus Terminal knows where every
command starts and stops. \keys{\mod+Shift+Up} and \keys{\mod+Shift+Down} jump from prompt to prompt,
\ui{Copy Last Command Output} grabs just what the last command printed, and a thin bar beside each prompt shows how it
went: green for success, red for failure. Hover the bar for the exit code and how long it took.

\shot{shell-integration.png}{\code{make test} failed (red bar), \code{ls} worked (green bar).}

For bash, add this to \code{\textasciitilde/.bashrc} on the machine the shell runs on (the user manual has the zsh
version):

\begin{verbatim}
PS0='\[\e]133;C\a\]'
PROMPT_COMMAND='printf "\033]133;D;%s\007" "$?"; '"${PROMPT_COMMAND}"
PS1='\[\e]133;A\a\]'"$PS1"'\[\e]133;B\a\]'
\end{verbatim}

\section{Pictures in the terminal}

Programs that print images the iTerm2 way (\code{imgcat}, \code{wezterm imgcat}, \code{chafa -f iterm}) show PNG,
JPEG and GIF images right in the scrollback. They scroll with the text and never spill past the edge.

\shot{inline-image.png}{An image printed by a script, with the text carrying on underneath.}

\section{Find}

\keys{\mod+F} opens the find bar. Toggle \ui{Aa} for case-sensitive and \ui{.*} for regular expressions.
\kbd{Enter} for the next match, \keys{Shift+Enter} for the previous one, \kbd{Esc} to close. It shows
``3 of 17'' so you know where you are.

\section{When the connection drops}

The pane dims and shows a \ui{Disconnected} card with \ui{Reconnect} and \ui{Close}. If auto-reconnect is on,
it says \ui{Reconnecting…} with a \ui{Cancel} button, and a toast tells you which attempt it's on.

If the problem was a password or host key, you get a friendly dialog explaining what happened, with
\ui{Retry}, \ui{Edit Profile}, \ui{Known Hosts} (for host-key trouble) and \ui{Close}.

% =========================================================================================
\chapter{Never type it twice}\label{ch:repeat}
\chapmeta{6 minutes}{You'll stop retyping the same commands, and let Nexus Terminal watch the logs for you}

\begin{tldr}
\textbf{Snippets} store commands. \textbf{Broadcast} types into many sessions at once. \textbf{Triggers} watch output and
tap you on the shoulder.
\end{tldr}

\section{Snippets: \keys{\mod+Shift+S}}

Save commands you use a lot, with a name, description and folder. Search them, click
\ui{Send to Session}, and the command is typed into the active session followed by \kbd{Enter}.

\crop{snippets.png}{633}{384}{611}{402}{0.8\linewidth}{The snippet library, organised into folders.}

Snippets can have blanks to fill in:

\begin{itemize}
  \item \code{\{\{hostname\}\}}, \code{\{\{date\}\}}, \code{\{\{time\}\}} and \code{\{\{timestamp\}\}} are filled in for you.
  \item Anything else, like \code{\{\{service\}\}}, is asked for when you send it.
\end{itemize}

So \code{journalctl -fu \{\{service\}\}} asks ``service?'' and tails whichever one you type.

\section{Broadcast: \keys{\mod+Shift+B}}

Turn broadcast on and everything you type goes to \emph{every} connected terminal. The status bar shows
\ui{BROADCAST} so you can't forget it's on.

\shot{broadcast.png}{Broadcast is on: the status bar says so in no uncertain terms.}

To broadcast to only some sessions, right-click the \iBroadcast{} button (or \menu{Session\then Broadcast
targets…}), tick the ones you want and click \ui{Apply \& Enable}.

\begin{danger}[With great power…]
Broadcast doesn't care which tab you \emph{meant}. Double-check the target list before you type anything
destructive. \code{sudo reboot} into twenty servers is very efficient and very memorable.
\end{danger}

\section{Triggers: let it watch for you}

\menu{Session\then Trigger Rules} (or the palette). Each rule has a pattern (plain text or a regular
expression), a case-sensitivity switch, an on/off switch, and an action:

\begin{center}\small
\begin{tabular}{@{}ll@{}}
  \ui{Notify} & shows a toast\\
  \ui{Highlight tab} & briefly outlines the pane\\
  \ui{Notify + Highlight} & both\\
  \ui{Sound} & asks your OS for attention (a bouncing dock icon or flashing taskbar button) and shows a toast\\
\end{tabular}
\end{center}

\crop{triggers.png}{725}{610}{725}{625}{0.75\linewidth}{Three rules: shout about errors, ping when a build finishes, and one switched off for now.}

Rules check each complete line of output, with colours stripped out first.

\begin{quickwin}
Add a rule for \code{BUILD SUCCESSFUL} with the \ui{Sound} action, start a long build, and go and make a cup of tea.
Nexus Terminal will tell you when it's done. (Or \code{(?i)error|fatal} for the opposite.)
\end{quickwin}

% =========================================================================================
\chapter{Files: SFTP and drag-and-drop}\label{ch:files}
\chapmeta{4 minutes}{You'll move files to and from servers without leaving the terminal}

\begin{tldr}
Click the \iFolder{} on an SSH tab. Or drop a file onto the terminal and choose \ui{Upload}.
\end{tldr}

\section{The SFTP browser}

Click the \iFolder{} folder icon on an SSH tab, \menu{Session\then SFTP browser}, or type ``sftp'' in the palette.
It opens in the side dock next to your terminal, in the folder your shell is in if the shell reports it
(fish does; bash and zsh need a one-line hook, see the user manual).

\shot{sftp.png}{Browsing the server's files next to the terminal.}

\begin{itemize}
  \item Move around with the breadcrumb trail, the \code{..} entry or \ui{Go to path}.
  \item Upload files; download one or many to a folder you choose.
  \item Rename, delete (it asks first), make folders, refresh.
  \item \iStar{} bookmarks the current folder, per profile.
  \item Ctrl-click (\ph{E1C4}-click) or Shift-click to select several things. Transfers show progress as they go.
\end{itemize}

\section{Drag and drop}

Drop files onto a terminal and choose:

\begin{description}
  \item[Upload to current directory] Pick the remote folder (SSH only).
  \item[Paste contents] Types the file's contents into the terminal, as if you'd typed it.
\end{description}

The upload picker looks like the SFTP browser but shows folders only, labelled \proto[amber]{READ ONLY} if you can't
write there and \proto[rose]{NO READ} if you can't look inside. Keyboard fans: \kbd{↑}\kbd{↓} to move, \kbd{Enter} to
choose, \kbd{Backspace} to go up, \kbd{Esc} to cancel. If the SFTP browser is open, the picker starts in the same folder.

% =========================================================================================
\chapter{Port forwards}\label{ch:forwards}
\chapmeta{3 minutes}{You'll reach a database or web app that only listens on the server}

\begin{tldr}
\emph{Local} forwards bring a remote port to your machine. \emph{Remote} forwards do the opposite.
\end{tldr}

Forwards saved in an SSH profile (under \ui{Port Forwarding}, with \ui{Active} ticked) start when you connect.
To manage them on a live session, right-click the terminal\then\ui{Port Forwards}, or \menu{Session\then Port forwards}.

\begin{description}
  \item[Local \iArrow] Listen on your computer, forward to somewhere the server can reach.
        Example: \code{127.0.0.1:15432}\then\code{127.0.0.1:5432} lets your desktop database app talk to Postgres on
        the server as if it were local.
  \item[Remote \ph{E058}] Ask the server to listen and send connections back to somewhere \emph{your} computer can reach.
        Great for showing a colleague the dev server running on your laptop.
\end{description}

Active forwards are listed with an \proto[accent]{L} or \proto[violet]{R} badge and a stop button. Toasts tell you when one
starts, stops or fails.

% =========================================================================================
\chapter{Recordings}\label{ch:recordings}
\chapmeta{2 minutes}{You'll have a replayable record of exactly what happened}

\keys{\mod+Shift+R} starts and stops recording the active session. A red dot appears on the tab and \ui{REC} in
the status bar. Recordings are standard asciicast v2 files, so other tools (like \code{asciinema}) can play them too.

\crop{recordings.png}{712}{392}{690}{415}{0.75\linewidth}{The recordings browser. Play, delete, or tidy up in bulk.}

The browser (sidebar footer, \menu{Session} menu or palette) shows how many recordings you have and how much space
they take. Play or delete each one, or use \ui{Clean up} to remove recordings older than 7, 30 or 90 days, or to keep
the total under 1\,GB or 5\,GB. Playback has play/pause, restart, a seek bar and speed control.

\begin{protip}
Record the next time you do something fiddly and rare (a database failover, a certificate rotation). Six months
later, ``how did we do that again?'' has an answer you can literally watch.
\end{protip}
